\subsection{Study Area Background}
\label{subsec:study_area}

The study area lies within the Choushui River Alluvial Fan (CRAF) on the western coastal plain of central Taiwan. Covering $\sim$2,000~km$^2$ across Changhua and Yunlin counties, the fan is bounded by the Wu River to the north, the Beigang River to the south, the Bagua Tableland and Douliu Hills to the east, and the Taiwan Strait to the west. Surface elevation decreases from approximately 100--150~m near the eastern highlands to only several meters above mean sea level near the coast \citep{survey_project_1999, liu_characterization_2004}.

Sediments beneath the CRAF were derived from the weathering and erosion of slate, quartzite, shale, sandstone, and mudstone in the upstream watershed \citep{liu_characterization_2004, hung2015_multiple}. Based on its hydrogeological characteristics, the fan is commonly divided into proximal, middle, and distal zones \citep{survey_project_1999}. Gravel and coarse sand dominate the proximal fan, where they form highly permeable deposits. Grain size generally decreases westward as fine sand, silt, and clay become increasingly interbedded across the middle and distal fans \citep{liu_characterization_2004, chen2001_stratigraphy}. This westward fining pattern produces a multilayer aquifer system in which permeable aquifers alternate with less permeable aquitards, as illustrated by the regional hydrogeological cross-section in \Cref{fig:studyarea}.

The CRAF is one of Taiwan's major agricultural regions and accounts for approximately 38\% of the country's rice production. Water demand rises sharply during the first rice cultivation season from mid-February to late March, within the dry season from November to April \citep{chang2022_wetanddry, chang_rice-field_2020}. To meet this seasonal demand, annual groundwater extraction has been estimated at 1.71 to 2.05 billion~m$^3$~yr$^{-1}$ \citep{tseng_estimating_2024}, with much of the extraction concentrated in the middle and distal fans \citep{hung_exploring_2024}. The associated decline in hydraulic head reduces pore-fluid pressure and transfers a greater share of the overburden load to the sediment skeleton, thereby increasing effective stress. Where this stress exceeds the preconsolidation stress of susceptible fine-grained aquitards and interbeds, the sediments undergo permanent compaction that contributes to land subsidence \citep{liu_characterization_2004, hung2015_multiple, gambolati_2015}. These coupled seasonal changes in groundwater levels and deformation make the middle and distal fans important settings for continuous monitoring.

\begin{figure}[p]
  \centering
  \begin{subfigure}[t]{0.50\textwidth}
    \centering
    \fbox{\parbox[c][1.20\linewidth][c]{0.96\linewidth}{\centering
      \placeholder{CRAF MAP WITH CHOUSHUI RIVER, MONITORING LOCATIONS, AND CROSS-SECTION TRACE}}}
    \caption{Location of the CRAF and the monitoring stations used in this study.}
    \label{fig:studyarea_a}
  \end{subfigure}

  \vspace{0.6\baselineskip}

  \begin{subfigure}[t]{\textwidth}
    \centering
    \fbox{\parbox[c][0.693\linewidth][c]{0.96\linewidth}{\centering
      \placeholder{REGIONAL HYDROGEOLOGICAL CROSS-SECTION ACROSS THE STUDY AREA}}}
    \caption{Regional hydrogeological cross-section across the CRAF.}
    \label{fig:studyarea_b}
  \end{subfigure}
  \caption{Study area and hydrogeological setting of the Choushui River Alluvial Fan.}
  \label{fig:studyarea}
\end{figure}

\FloatBarrier
